{\def\LTcaptype{none} % do not increment counter
\begin{longtable}[]{@{}
  >{\raggedright\arraybackslash}p{(\linewidth - 2\tabcolsep) * \real{0.5000}}
  >{\raggedright\arraybackslash}p{(\linewidth - 2\tabcolsep) * \real{0.5000}}@{}}
\toprule\noalign{}
\endhead
\bottomrule\noalign{}
\endlastfoot
\textbf{Course} & Cognitive Modeling (elective specialization course) \\
\textbf{Credits} & 2 units: 1 theory unit (16 h) + 1 practical unit (32
h) = 48 h \\
\textbf{Format} & 16 weeks × (1 h lecture + 2 h hands-on lab) \\
\textbf{Audience} & MSc students and final-year BSc students in
cognitive science, psychology, computer science and AI \\
\textbf{Prerequisites} & Statistical methods in cognitive science;
experimental methods in cognitive psychology. Basic probability,
statistics and introductory programming are assumed; Bayesian statistics
is \emph{not} assumed. \\
\textbf{Instructor} & \emph{{[}name, office, e-mail, office hours{]}} \\
\textbf{Teaching assistant} & \emph{{[}name, e-mail, lab hours{]}} \\
\textbf{Course page} & \emph{{[}LMS link{]}} \\
\end{longtable}
}

\begin{center}\rule{0.5\linewidth}{0.5pt}\end{center}

\subsection{Course description}

Cognitive models are explicit, mathematical or computational
descriptions of the processes behind behavior. This course teaches how
to build such models, fit them to experimental data, evaluate them, and
interpret what they say about the mind. We move from intuition to
formalism to code to interpretation: from the logic of a model, through
least-squares and maximum-likelihood estimation, model evaluation and
Bayesian inference, to Markov chain Monte Carlo and hierarchical models.
We finish with models from memory, perception, categorization, decision
making and learning, and their links to artificial intelligence. Every
week includes a two-hour lab in Python in which students fit, evaluate
and interpret a model on (simulated or real) experimental data.

\subsection{Learning goals}

\textbf{General goals.} By the end of the course students will be
familiar with: 1. different approaches to and models of cognitive
modeling; 2. how to design a model to address a cognitive problem; 3.
applications of modeling in cognitive science and artificial
intelligence.

\textbf{Specific goals.} Students will gain the skills to: 1. implement
modeling methods on experimental data; 2. interpret research results
obtained with cognitive models.

\subsection{Topics}

\begin{enumerate}
\def\labelenumi{\arabic{enumi}.}
\tightlist
\item
  Fundamentals of cognitive modeling
\item
  Parameter estimation methods (least squares / mean squared error,
  maximum likelihood)
\item
  Evaluation of cognitive models
\item
  Bayesian modeling
\item
  Parameter estimation with data-driven methods (Monte Carlo,
  Metropolis--Hastings and Gibbs sampling, MCMC)
\item
  Examples of modeling in different areas of cognitive science
\end{enumerate}

\subsection{Weekly schedule}

{\def\LTcaptype{none} % do not increment counter
\begin{longtable}[]{@{}
  >{\raggedright\arraybackslash}p{(\linewidth - 8\tabcolsep) * \real{0.2000}}
  >{\raggedright\arraybackslash}p{(\linewidth - 8\tabcolsep) * \real{0.2000}}
  >{\raggedright\arraybackslash}p{(\linewidth - 8\tabcolsep) * \real{0.2000}}
  >{\raggedright\arraybackslash}p{(\linewidth - 8\tabcolsep) * \real{0.2000}}
  >{\raggedright\arraybackslash}p{(\linewidth - 8\tabcolsep) * \real{0.2000}}@{}}
\toprule\noalign{}
\begin{minipage}[b]{\linewidth}\raggedright
Week
\end{minipage} & \begin{minipage}[b]{\linewidth}\raggedright
Chapter
\end{minipage} & \begin{minipage}[b]{\linewidth}\raggedright
Lecture
\end{minipage} & \begin{minipage}[b]{\linewidth}\raggedright
Lab
\end{minipage} & \begin{minipage}[b]{\linewidth}\raggedright
Due
\end{minipage} \\
\midrule\noalign{}
\endhead
\bottomrule\noalign{}
\endlastfoot
1 & 1 & What is a cognitive model? Marr's levels; notation & Install the
course environment; simulate a random walk & --- \\
2 & 1 & Anatomy of a model: SDT and forgetting curves & SDT, ROC/zROC;
forgetting curves with binomial noise & A0 \\
3 & 2 & Least squares, error surfaces and optimization & Grid search,
Nelder--Mead, L-BFGS-B; failure modes & --- \\
4 & 2 & Maximum likelihood & Binomial MLE; \(G^2\); multi-condition SDT
& A1 \\
5 & 2 & Uncertainty; fitting many participants & Hessian, profile and
bootstrap CIs; aggregation artifacts & --- \\
6 & 3 & Fit vs.~generalization; cross-validation & Overfitting; K-fold
CV; generalization criterion & A2 \\
7 & 3 & LRT, AIC/BIC; parameter and model recovery & Recovery studies;
confusion matrices & Project teams \\
8 & 1--3 & Midterm review: the whole pipeline & \textbf{Midterm exam}
(90 min) + project consultation & --- \\
9 & 4 & Bayesian modeling I: priors, posteriors, grids & Beta-binomial;
2-D grid posterior for SDT; prior predictive checks & A3, project
proposal \\
10 & 4 & Bayesian modeling II: PPCs and Bayes factors & PPC; Bayes
factors by grid and Savage--Dickey; prior sensitivity & --- \\
11 & 5 & Monte Carlo and Metropolis--Hastings & Metropolis from scratch;
tuning & --- \\
12 & 5 & Gibbs sampling; convergence diagnostics & Gibbs from scratch;
R-hat and ESS; ArviZ & A4 \\
13 & 5 & Probabilistic programming; hierarchical models & Hierarchical
SDT in PyMC & Progress check \\
14 & 6 & Categorization (GCM) and decisions (drift-diffusion) & GCM
vs.~prototype; EZ-diffusion & A5 \\
15 & 6 & Reinforcement learning, Bayesian cognition and AI & Q-learning
fits and recovery; ideal observer & --- \\
16 & --- & Synthesis and good modeling practice & \textbf{Project
presentations} & Presentations \\
Exam period & 1--6 & --- & \textbf{Final exam} (120 min) & Report + code
(before the exam) \\
\end{longtable}
}

Detailed materials for each week are in \nolinkurl{weekNN_*.md}.

\subsection{Assessment (0--20 scale)}

{\def\LTcaptype{none} % do not increment counter
\begin{longtable}[]{@{}
  >{\raggedright\arraybackslash}p{(\linewidth - 2\tabcolsep) * \real{0.5000}}
  >{\raggedright\arraybackslash}p{(\linewidth - 2\tabcolsep) * \real{0.5000}}@{}}
\toprule\noalign{}
\begin{minipage}[b]{\linewidth}\raggedright
Component
\end{minipage} & \begin{minipage}[b]{\linewidth}\raggedright
Points
\end{minipage} \\
\midrule\noalign{}
\endhead
\bottomrule\noalign{}
\endlastfoot
Lab check-offs (weekly in-lab exercises) & 2.0 \\
Assignments A1--A5 (1.0 each; A0 required, ungraded) & 5.0 \\
In-lecture quizzes (about 10; lowest 2 dropped) & 1.0 \\
Midterm exam (Week 8) & 3.0 \\
Final project: proposal 1.0 · presentation 2.0 · report + code 3.0 &
6.0 \\
Final exam (cumulative, emphasis on Chapters 4--6) & 3.0 \\
\textbf{Total} & \textbf{20.0} \\
\end{longtable}
}

Details, rubrics and policies are in \nolinkurl{assessments/README.md}.

\subsection{Policies (summary)}

\begin{itemize}
\tightlist
\item
  \textbf{Late work:} −10\% per day, maximum 3 days.
\item
  \textbf{Collaboration:} discuss freely; write your own code and text
  (the project is a team product for pairs). Acknowledge discussions and
  reused code.
\item
  \textbf{AI tools:} allowed for debugging, explanations and language
  editing. State how you used them; you must be able to explain
  everything you submit.
\item
  \textbf{Exams:} closed book, formula sheet provided, non-programmable
  calculator allowed.
\item
  \textbf{Attendance:} labs are where most learning happens; lab
  check-offs require attendance (or an agreed alternative).
\item
  \textbf{Accessibility:} students needing accommodations should contact
  the instructor in the first two weeks.
\end{itemize}

\subsection{Software}

Python 3.11+ with numpy, scipy, pandas, matplotlib, JupyterLab, PyMC and
ArviZ, installed on students' own laptops in Week 1
(Miniforge/conda-forge recommended; Google Colab as fallback). See
\nolinkurl{environment.yml}, \nolinkurl{requirements.txt} and the Week 1
lab.

\subsection{Readings}

Main references (see \href{reading_list.md}{\nolinkurl{reading_list.md}}
for the chapter mapping): - Farrell, S., \& Lewandowsky, S. (2018).
\emph{Computational modeling of cognition and behavior.} Cambridge
University Press. \textbf{(primary text, Weeks 1--8)} - Lee, M. D., \&
Wagenmakers, E.-J. (2014). \emph{Bayesian cognitive modeling: A
practical course.} Cambridge University Press. \textbf{(primary text,
Weeks 9--14)} - Kruschke, J. K. (2014). \emph{Doing Bayesian data
analysis: A tutorial with R, JAGS, and Stan} (2nd ed.). Academic Press.
- Martin, O. A., Kumar, R., \& Lao, J. (2021). \emph{Bayesian modeling
and computation in Python.} Chapman and Hall/CRC. - Anderson, B. (2014).
\emph{Computational neuroscience and cognitive modelling: A student's
introduction to methods and procedures.} SAGE. - Verguts, T. (2022).
\emph{Introduction to modeling cognitive processes.} MIT Press. -
Ghodousian, A. (1401). \emph{Mabani-ye ravesh-haye mostaghim dar
behine-sazi gheyr-khatti} {[}Foundations of direct methods in nonlinear
optimization{]} (in Persian). - Azimifar, Z., \& Ahmadi, E. (1396).
\emph{Shenasayi-ye amari-ye olgu: yadgiri-ye nazari-ye moghaddamati}
{[}Statistical pattern recognition: Introductory learning theory{]},
vol.~1 (in Persian).
